% ============================================================
% Part 2: Probability & Statistics, Information Theory,
%         Logic & Computation (Slides 26-48)
% Companion lecture: "The Building Blocks of LLMs"
% This file is \input'ed from the main document.
% ============================================================


% ************************************************************
% SECTION 3: PROBABILITY & STATISTICS (Slides 26-34)
% "The Language of Uncertainty"
% Section color: aiViolet
% ************************************************************
\section{Probability \& Statistics}

% ----------------------------------------------------------
% Slide 26: Section Divider
% ----------------------------------------------------------
{
\setbeamercolor{background canvas}{bg=aiViolet}
\begin{frame}[plain]
  \centering
  \vspace{2cm}
  \textcolor{white}{\Huge\bfseries Probability \& Statistics}\\[0.5em]
  \textcolor{white}{\Huge\bfseries The Language of Uncertainty}\\[1em]
  \textcolor{white!80}{\large 1654--1933}\\[1.5em]
  \textcolor{white!60}{\normalsize Pascal \(\cdot\) Fermat \(\cdot\) Bernoulli \(\cdot\) Bayes \(\cdot\) Kolmogorov \(\cdot\) Fisher \(\cdot\) Shannon \(\cdot\) Boltzmann}

  \note{
    Transition from Section~2: ``The network can now learn. But learn WHAT?
    It learns to predict probability distributions over words.
    For that, we need probability theory.''

    Tagline: An LLM doesn't KNOW the next word. It assigns PROBABILITIES
    to every possible word.

    This section covers \(\sim\)10 minutes.
  }
\end{frame}
}

% ----------------------------------------------------------
% Slide 27: The Birth of Probability -- A Gambling Problem
% ----------------------------------------------------------
\begin{frame}{The Birth of Probability --- A Gambling Problem (1654)}

  \begin{columns}[T]
    \begin{column}{0.48\textwidth}
      \begin{personbox}[Blaise Pascal (1623--1662)]{aiViolet}
        \textbf{Origin:} Clermont-Ferrand, France\\[3pt]
        Mathematician, physicist, philosopher. Co-founded probability
        theory in correspondence with Fermat (1654) about the
        \alert{problem of points}: how to divide stakes in an
        interrupted game of chance.
      \end{personbox}

      \vspace{0.4em}
      \begin{personbox}[Pierre de Fermat (1601--1665)]{aiViolet}
        \textbf{Origin:} Beaumont-de-Lomagne, France\\[3pt]
        His reply to Pascal laid out the counting arguments that
        became \textbf{combinatorial probability}.
      \end{personbox}
    \end{column}

    \begin{column}{0.49\textwidth}
      % Visual: Letter exchange and game tree
      \begin{tikzpicture}[
        letter/.style={draw=aiViolet!60, rounded corners, fill=aiViolet!5,
          text width=2cm, align=center, font=\scriptsize, inner sep=4pt},
        >=Stealth
      ]
        \node[letter] (pascal) at (0,2.5) {Pascal\\Paris};
        \node[letter] (fermat) at (3.5,2.5) {Fermat\\Toulouse};
        \draw[->,thick,aiViolet!70] (pascal.east) -- node[above,font=\tiny] {1654} (fermat.west);
        \draw[<-,thick,aiViolet!70] ([yshift=-3pt]pascal.east) -- ([yshift=-3pt]fermat.west);

        % Simplified game tree
        \node[font=\footnotesize\bfseries,aiViolet] at (1.75,1.4) {Problem of Points};
        \node[circle,draw=aiViolet,fill=aiViolet!10,inner sep=1pt,font=\tiny] (r) at (1.75,0.8) {Game};
        \node[circle,draw=aiViolet!60,fill=aiViolet!5,inner sep=1pt,font=\tiny] (a) at (0.5,0) {A wins};
        \node[circle,draw=aiViolet!60,fill=aiViolet!5,inner sep=1pt,font=\tiny] (b) at (3,0) {B wins};
        \draw[->,aiViolet!60] (r) -- (a);
        \draw[->,aiViolet!60] (r) -- (b);
        \node[font=\tiny\itshape,text=gray,text width=4cm,align=center] at (1.75,-0.6)
          {How to split the pot if the game stops early?};
      \end{tikzpicture}
    \end{column}
  \end{columns}

  \note{
    ``Probability theory was born because a gambler asked Pascal a question
    about how to split the pot in an unfinished dice game. Pascal wrote to
    Fermat. Their letters invented a new branch of mathematics.''

    Fermat again! We saw him in calculus (tangent lines, Slide~19). He
    contributed foundational ideas to TWO separate pillars of the LLM.
  }
\end{frame}

% ----------------------------------------------------------
% Slide 28: From Gambling to Laws -- Bernoulli and Laplace
% ----------------------------------------------------------
\begin{frame}{From Gambling to Laws --- Bernoulli and Laplace}

  \begin{columns}[T]
    \begin{column}{0.5\textwidth}
      \begin{personbox}[Jacob Bernoulli (1655--1705)]{aiViolet}
        \textbf{Origin:} Basel, Switzerland\\[3pt]
        Proved the \alert{law of large numbers} (published posthumously
        in \emph{Ars Conjectandi}, 1713): as you repeat an experiment,
        the observed frequency converges to the true probability.
      \end{personbox}

      \vspace{0.3em}
      \begin{personbox}[Pierre-Simon Laplace (1749--1827)]{aiViolet}
        \textbf{Origin:} Beaumont-en-Auge, France\\[3pt]
        \emph{Th\'eorie analytique des probabilit\'es} (1812)
        systematized the entire field. Popularized Bayes' theorem.
      \end{personbox}
    \end{column}

    \begin{column}{0.47\textwidth}
      % Law of large numbers illustration
      \begin{tikzpicture}[scale=0.85]
        \begin{axis}[
          width=6cm, height=4cm,
          xlabel={\scriptsize Number of flips},
          ylabel={\scriptsize Running average},
          ymin=0.3, ymax=0.7,
          xmin=0, xmax=100,
          ytick={0.3,0.4,0.5,0.6,0.7},
          xtick={0,25,50,75,100},
          grid=major,
          grid style={gray!20},
          axis lines=left,
          tick label style={font=\tiny},
          label style={font=\tiny},
        ]
          % Simulated convergence to 0.5
          \addplot[aiViolet, thick, smooth] coordinates {
            (1,1.0) (3,0.67) (5,0.6) (10,0.6) (15,0.53)
            (20,0.55) (30,0.47) (40,0.52) (50,0.48)
            (60,0.51) (70,0.49) (80,0.50) (100,0.50)
          };
          \draw[dashed, accentOrange, thick] (axis cs:0,0.5) -- (axis cs:100,0.5);
          \node[font=\tiny, accentOrange] at (axis cs:85,0.54) {$p = 0.5$};
        \end{axis}
      \end{tikzpicture}

      \vspace{0.3em}
      \begin{insightbox}
        Why do LLMs improve with more data?
        Bernoulli's law: more data \(\Rightarrow\) learned probabilities
        converge to the true patterns of language.
      \end{insightbox}
    \end{column}
  \end{columns}

  \note{
    ``Flip a coin 10 times: anything could happen.
    Flip it 10 million times: the average WILL be 50\%.''

    LLM connection: LLMs get better with more training data because the
    law of large numbers ensures that the learned word-frequency
    probabilities converge to the true distribution.
  }
\end{frame}

% ----------------------------------------------------------
% Slide 29: Bayes' Theorem -- Updating Beliefs with Evidence
% ----------------------------------------------------------
\begin{frame}{Bayes' Theorem --- Updating Beliefs with Evidence (1763)}

  \begin{columns}[T]
    \begin{column}{0.48\textwidth}
      \begin{personbox}[{Thomas Bayes (c.\,1701--1761)}]{aiViolet}
        \textbf{Origin:} London, England\\[3pt]
        Presbyterian minister and mathematician.
        \emph{Essay towards solving a Problem in the Doctrine of Chances}
        published \textbf{posthumously} (1763) by his friend Richard
        Price.
      \end{personbox}

      \vspace{0.4em}
      % Conceptual recipe
      \begin{tcolorbox}[colback=aiViolet!5, colframe=aiViolet!60,
        fonttitle=\bfseries\small, title={Recipe: Updating Beliefs},
        rounded corners, boxrule=0.5pt, left=3pt, right=3pt, top=2pt, bottom=2pt]
        \footnotesize
        New belief =\\[2pt]
        \(\dfrac{\text{(how likely the evidence if my guess is right)}
          \;\times\; \text{(original guess)}}
          {\text{(how likely the evidence overall)}}\)
      \end{tcolorbox}
    \end{column}

    \begin{column}{0.49\textwidth}
      % Formal inset
      \begin{tcolorbox}[colback=gray!5, colframe=gray!40,
        fonttitle=\bfseries\small, title={Formal},
        rounded corners, boxrule=0.5pt, left=3pt, right=3pt, top=2pt, bottom=2pt]
        \[
          P(A \mid B) = \frac{P(B \mid A)\; P(A)}{P(B)}
        \]
      \end{tcolorbox}

      \vspace{0.3em}
      % Surprising example
      \begin{tcolorbox}[colback=accentOrange!5, colframe=accentOrange!60,
        fonttitle=\bfseries\small, title={Surprise Example},
        rounded corners, boxrule=0.5pt, left=3pt, right=3pt, top=2pt, bottom=2pt]
        \footnotesize
        Test accuracy: 99\%\\
        Disease prevalence: 1 in 10,000\\
        Positive test \(\Rightarrow\) chance you have it?\\[3pt]
        \textbf{Answer: \(\sim\)1\%}\\
        Most positives are \emph{false} positives!
      \end{tcolorbox}

      \vspace{0.3em}
      \begin{insightbox}
        LLMs implicitly learn conditional probabilities:\\
        \(P(\text{next word} \mid \text{previous words})\)
      \end{insightbox}
    \end{column}
  \end{columns}

  \note{
    ``Bayes was a minister, not a professor. He published almost nothing
    in his lifetime. His one big idea --- how to update beliefs from
    evidence --- was published after his death by a friend. It became one
    of the most important ideas in all of science.''

    LLM connection: While LLMs don't explicitly compute Bayes' theorem,
    training causes them to learn conditional probabilities. That IS
    Bayesian reasoning, learned from data.

    \verifytag{Bayes birth year uncertain --- some sources give 1701,
    others 1702.}
  }
\end{frame}

% ----------------------------------------------------------
% Slide 30: Kolmogorov's Axioms [EXTENDED]
% ----------------------------------------------------------
\begin{frame}{Kolmogorov's Axioms --- Probability Made Rigorous (1933)}

  \begin{columns}[T]
    \begin{column}{0.45\textwidth}
      \begin{personbox}[Andrey Kolmogorov (1903--1987)]{aiViolet}
        \textbf{Origin:} Tambov, Russia\\[3pt]
        Published \emph{Grundbegriffe der Wahrscheinlichkeitsrechnung}
        (Foundations of the Theory of Probability, 1933), axiomatizing
        probability using measure theory.
      \end{personbox}

      \vspace{0.3em}
      {\small\itshape This put probability on the same rigorous footing
      as geometry (Euclid) and calculus (Cauchy/Weierstrass).}

      \vspace{0.3em}
      {\footnotesize\textcolor{gray}{[EXTENDED] --- can be cut without
      losing the narrative}}
    \end{column}

    \begin{column}{0.52\textwidth}
      % Three axioms in plain English
      \begin{tcolorbox}[colback=aiViolet!5, colframe=aiViolet,
        fonttitle=\bfseries, title={Kolmogorov's Three Axioms},
        rounded corners, boxrule=1pt, left=4pt, right=4pt, top=3pt, bottom=3pt]
        \begin{enumerate}\setlength\itemsep{6pt}
          \item \textbf{Non-negativity:} Probabilities are between 0 and~1.
          \item \textbf{Certainty:} Something must happen\\
                (total probability = 1).
          \item \textbf{Additivity:} Probabilities of mutually exclusive
                events add up.
        \end{enumerate}
      \end{tcolorbox}

      \vspace{0.4em}
      {\small\centering Three simple rules. Every probability calculation
      in every LLM obeys them.\par}
    \end{column}
  \end{columns}

  \note{
    This slide is marked [EXTENDED] and can be cut if time is short.

    The conceptual point: before Kolmogorov, probability was a collection
    of tricks. After Kolmogorov, it was a branch of mathematics as
    rigorous as geometry.

    \verifytag{Publication year 1933 and German title confirmed in
    standard references.}
  }
\end{frame}

% ----------------------------------------------------------
% Slide 31: Maximum Likelihood -- Finding the Best Fit
% ----------------------------------------------------------
\begin{frame}{Maximum Likelihood --- Finding the Best Fit}

  \begin{columns}[T]
    \begin{column}{0.47\textwidth}
      \begin{personbox}[Ronald A. Fisher (1890--1962)]{aiViolet}
        \textbf{Origin:} London, England\\[3pt]
        Introduced \alert{maximum likelihood estimation} (concept 1912,
        formalized 1922). His 1922 paper ``On the Mathematical
        Foundations of Theoretical Statistics'' established MLE as a
        foundational method.
      \end{personbox}

      \vspace{0.3em}
      % Conceptual recipe
      \begin{tcolorbox}[colback=aiViolet!5, colframe=aiViolet!60,
        fonttitle=\bfseries\small, title={Recipe: Maximum Likelihood},
        rounded corners, boxrule=0.5pt, left=3pt, right=3pt, top=2pt, bottom=2pt]
        \footnotesize
        Given the data I observed, which parameters make
        this data \textbf{most probable}?\\[2pt]
        \(\rightarrow\) Choose the best fit.
      \end{tcolorbox}
    \end{column}

    \begin{column}{0.5\textwidth}
      % Visual: curve fitting
      \begin{tikzpicture}[scale=0.85]
        \begin{axis}[
          width=6cm, height=4.5cm,
          axis lines=left,
          xlabel={\scriptsize $x$},
          ylabel={\scriptsize Probability},
          xtick=\empty, ytick=\empty,
          xmin=-3, xmax=3,
          ymin=0, ymax=0.5,
          tick label style={font=\tiny},
          label style={font=\tiny},
        ]
          % Data histogram hint (points)
          \addplot[only marks, mark=*, mark size=1.5pt, aiViolet]
            coordinates {(-2,0.06) (-1.5,0.13) (-1,0.22) (-0.5,0.33)
              (0,0.40) (0.5,0.35) (1,0.24) (1.5,0.12) (2,0.05)};
          % Candidate curves (light)
          \addplot[gray!40, thick, smooth, domain=-3:3, samples=40]
            {0.4*exp(-x*x/0.5)};
          \addplot[gray!40, thick, smooth, domain=-3:3, samples=40]
            {0.3*exp(-(x-0.5)*(x-0.5)/2)};
          % MLE curve (bold)
          \addplot[aiViolet, very thick, smooth, domain=-3:3, samples=40]
            {0.4*exp(-x*x/2)};
          \node[font=\tiny\bfseries, aiViolet] at (axis cs:1.8,0.35)
            {MLE fit};
        \end{axis}
      \end{tikzpicture}

      \vspace{0.2em}
      \begin{insightbox}
        LLM training is MLE at scale: find weights that maximize
        the probability of the training text. Fisher's method from 1922,
        applied to trillions of words.
      \end{insightbox}
    \end{column}
  \end{columns}

  \note{
    ``Fisher introduced the concept in 1912 and formalized it in 1922.
    He was also a pioneering statistician and controversial eugenicist ---
    we mention the math, not the politics.''

    LLM connection: LLM training is essentially maximum likelihood
    estimation: find the model parameters (weights) that maximize the
    probability of the training text.

    \verifytag{Fisher introduced concept in 1912 undergraduate paper,
    formalized in 1922 --- confirm both dates.}
  }
\end{frame}

% ----------------------------------------------------------
% Slide 32: Statistical Language Models -- Shannon's Breakthrough
% ----------------------------------------------------------
\begin{frame}{Statistical Language Models --- Shannon's Breakthrough (1948)}

  \begin{columns}[T]
    \begin{column}{0.45\textwidth}
      \begin{personbox}[Claude Shannon (1916--2001)]{aiViolet}
        \textbf{Origin:} Petoskey, Michigan, USA\\[3pt]
        Published \emph{A Mathematical Theory of Communication} (1948),
        founding information theory. Modeled English as a stochastic
        process and built \textbf{n-gram language models} decades before
        AI.
      \end{personbox}

      \vspace{0.3em}
      % Diagram: Shannon to LLM
      \begin{tikzpicture}[
        box/.style={draw=aiViolet!60, rounded corners, fill=aiViolet!5,
          font=\scriptsize, inner sep=3pt, text width=2.8cm, align=center},
        >=Stealth
      ]
        \node[box] (sh) at (0,0) {Shannon (1948):\\predict next letter\\using previous $N$ letters};
        \node[box] (llm) at (0,-1.6) {LLM (2024):\\predict next token\\using ALL previous tokens};
        \draw[->,thick,aiViolet!70] (sh) -- (llm);
      \end{tikzpicture}
    \end{column}

    \begin{column}{0.52\textwidth}
      % Shannon's text examples
      \begin{tcolorbox}[colback=gray!5, colframe=gray!50,
        fonttitle=\bfseries\small,
        title={Shannon's Approximations to English},
        rounded corners, boxrule=0.5pt, left=3pt, right=3pt, top=2pt, bottom=2pt]
        \scriptsize
        \textbf{0th order} (random letters):\\
        {\ttfamily XFOML RXKHRJFFJUJ}\\[2pt]
        \textbf{1st order} (letter frequencies):\\
        {\ttfamily OCRO HLI RGWR NMIE}\\[2pt]
        \textbf{2nd order} (letter pairs):\\
        {\ttfamily ON IE ANTSOUTINYS}\\[2pt]
        \textbf{Word-level bigram}:\\
        {\ttfamily THE HEAD AND IN FRONTAL ATTACK}\\[3pt]
        \textit{More context \(\Rightarrow\) more readable!}
      \end{tcolorbox}

      \vspace{0.2em}
      \begin{insightbox}
        Shannon was doing language modeling in 1948 --- 70~years
        before GPT. The same core idea, vastly more data.
      \end{insightbox}
    \end{column}
  \end{columns}

  \note{
    ``Shannon showed that language has statistical structure, and that you
    can exploit that structure to predict what comes next. That's exactly
    what every LLM does.''

    Shannon will appear AGAIN in Section 4 (Information Theory, Slide~36)
    for his entropy formula. Flag to students: ``We'll see Shannon again
    very soon --- his 1948 paper is so foundational it contributes to
    multiple sections.''
  }
\end{frame}

% ----------------------------------------------------------
% Slide 33: The Softmax Function -- From Thermodynamics to
%           Token Prediction
% ----------------------------------------------------------
\begin{frame}{The Softmax Function --- From Physics to Token Prediction}

  \begin{columns}[T]
    \begin{column}{0.46\textwidth}
      \begin{personbox}[Ludwig Boltzmann (1844--1906)]{aiViolet}
        \textbf{Origin:} Vienna, Austria\\[3pt]
        His \alert{Boltzmann distribution} (1870s) describes the
        probability of a physical system being in a state based on its
        energy: \(P \propto e^{-E/kT}\).
      \end{personbox}

      \vspace{0.3em}
      {\footnotesize
      \textbf{Lineage:}
      Boltzmann (1870s, physics) \(\to\)
      Gibbs (1900s, statistical mechanics) \(\to\)
      Luce (1959, psychology of choice) \(\to\)
      \textbf{Softmax} (ML)}

      \vspace{0.3em}
      % Conceptual recipe
      \begin{tcolorbox}[colback=aiViolet!5, colframe=aiViolet!60,
        fonttitle=\bfseries\small, title={Recipe: Softmax},
        rounded corners, boxrule=0.5pt, left=3pt, right=3pt, top=2pt, bottom=2pt]
        \footnotesize
        1. Take each raw score, raise $e$ to that power\\
        {\scriptsize\itshape (makes everything positive, magnifies differences)}\\
        2. Add up all the results\\
        3. Divide each by the total\\[2pt]
        \(\Rightarrow\) Now they are probabilities that sum to~1!
      \end{tcolorbox}
    \end{column}

    \begin{column}{0.51\textwidth}
      % Softmax pipeline
      \begin{tikzpicture}[
        softstep/.style={draw=aiViolet!60, rounded corners, fill=aiViolet!5,
          font=\scriptsize, inner sep=3pt, text width=4.2cm, align=left},
        >=Stealth
      ]
        \node[softstep] (s1) at (0, 3.0)
          {\textbf{Logits:} $[2.1,\; 0.8,\; 0.3,\; {-}1.5]$};
        \node[softstep] (s2) at (0, 2.0)
          {\textbf{Exp:} $[8.17,\; 2.23,\; 1.35,\; 0.22]$};
        \node[softstep] (s3) at (0, 1.0)
          {\textbf{Normalize:} $[0.68,\; 0.19,\; 0.11,\; 0.02]$};
        \draw[->,aiViolet!70] (s1.south) -- (s2.north);
        \draw[->,aiViolet!70] (s2.south) -- (s3.north);

        % Bar chart
        \fill[aiViolet!70] (0, -0.1) rectangle (2.3, 0.3);
        \fill[aiViolet!50] (2.5, -0.1) rectangle (3.15, 0.3);
        \fill[aiViolet!35] (3.35, -0.1) rectangle (3.72, 0.3);
        \fill[aiViolet!20] (3.92, -0.1) rectangle (3.99, 0.3);
        \node[font=\tiny,below] at (1.15, -0.15) {``the'' 68\%};
        \node[font=\tiny,below] at (2.82, -0.15) {``a'' 19\%};
        \node[font=\tiny,below] at (3.53, -0.15) {``that''};
        \node[font=\tiny,below] at (3.95, -0.15) {\ldots};
      \end{tikzpicture}

      \vspace{0.3em}
      % Formal inset
      \begin{tcolorbox}[colback=gray!5, colframe=gray!40,
        fonttitle=\bfseries\small, title={Formal},
        rounded corners, boxrule=0.5pt, left=3pt, right=3pt, top=1pt, bottom=1pt]
        \footnotesize
        $\displaystyle \mathrm{softmax}(z_i) = \frac{e^{z_i}}{\sum_j e^{z_j}}$
        \hfill {\scriptsize $\sum$ = ``add up for all $j$''}
      \end{tcolorbox}
    \end{column}
  \end{columns}

  \note{
    ``Boltzmann invented this formula to describe how gas molecules
    distribute among energy states. 150 years later, the exact same
    formula tells an LLM which word to say next.''

    INTERACTIVE MOMENT: Quick poll --- ``What word comes next:
    `The cat sat on the \rule{1cm}{0.15mm}'?'' Collect student answers,
    then show the AI's softmax distribution for the same prompt.

    \verifytag{Boltzmann distribution typically dated to 1870s; ``1868''
    may refer to early papers. R.~Duncan Luce's 1959 book
    \emph{Individual Choice Behavior} --- confirm.}
  }
\end{frame}

% ----------------------------------------------------------
% Slide 34: Convergence Diagram Update -- Probability
% ----------------------------------------------------------
\begin{frame}{Building the Transformer --- Probability \& Statistics}

  \begin{columns}[T]
    \begin{column}{0.45\textwidth}
      % Convergence diagram placeholder
      \visualplaceholder[5cm]{Convergence diagram: Sections 1--2 lit
        (Linear Algebra in blue, Calculus in green). Section~3
        ``Probability \& Statistics'' now lights up in orange/violet.
        Label: ``Softmax = next-token prediction.''}
    \end{column}

    \begin{column}{0.52\textwidth}
      \begin{tcolorbox}[colback=aiViolet!5, colframe=aiViolet,
        fonttitle=\bfseries, title={Section 3 Summary},
        rounded corners, boxrule=1pt, left=4pt, right=4pt, top=3pt, bottom=3pt]
        \footnotesize
        \begin{tabular}{@{}ll@{}}
          \toprule
          \textbf{Who / When} & \textbf{Contribution}\\
          \midrule
          Pascal / Fermat (1654) & Birth of probability\\
          Bernoulli (1713) & Law of large numbers\\
          Bayes (1763) / Laplace & Updating beliefs from evidence\\
          Kolmogorov (1933) & Rigorous axioms\\
          Fisher (1922) & Maximum likelihood estimation\\
          Shannon (1948) & Language as a statistical process\\
          Boltzmann (1870s) & \(\to\) Luce (1959) \(\to\) Softmax\\
          \bottomrule
        \end{tabular}
      \end{tcolorbox}

      \vspace{0.3em}
      \begin{insightbox}
        Every time an LLM generates a word, it computes a softmax
        probability distribution. The math started with a gambling
        question in 1654.
      \end{insightbox}
    \end{column}
  \end{columns}

  \note{
    Connection arrow: ``Every time an LLM generates a word, it computes
    a softmax probability distribution. The math started with a gambling
    question in 1654.''

    Transition: ``The network predicts probabilities. But how do we
    measure whether those probabilities are GOOD? For that, we need
    information theory.''
  }
\end{frame}


% ************************************************************
% SECTION 4: INFORMATION THEORY (Slides 35-41)
% "Measuring Meaning"
% Section color: aiViolet
% ************************************************************
\section{Information Theory}

% ----------------------------------------------------------
% Slide 35: Section Divider
% ----------------------------------------------------------
{
\setbeamercolor{background canvas}{bg=aiViolet}
\begin{frame}[plain]
  \centering
  \vspace{2cm}
  \textcolor{white}{\Huge\bfseries Information Theory}\\[0.5em]
  \textcolor{white}{\Huge\bfseries Measuring Meaning}\\[1em]
  \textcolor{white!80}{\large 1865--1948}\\[1.5em]
  \textcolor{white!60}{\normalsize Clausius \(\cdot\) Boltzmann \(\cdot\) Shannon \(\cdot\) Kullback \(\cdot\) Zipf}

  \note{
    Transition: ``The network predicts probabilities. But how do we
    measure whether those probabilities are GOOD? For that, we need
    information theory.''

    Tagline: How surprised should you be by the next word? Information
    theory measures that.

    This section covers \(\sim\)8 minutes.
  }
\end{frame}
}

% ----------------------------------------------------------
% Slide 36: Entropy -- From Steam Engines to Bits
% ----------------------------------------------------------
\begin{frame}{Entropy --- Invented Three Times}

  \begin{columns}[T]
    \begin{column}{0.32\textwidth}
      % Column 1: Clausius
      \begin{tcolorbox}[colback=aiViolet!5, colframe=aiViolet!60,
        fonttitle=\bfseries\scriptsize, title={Clausius (1865)},
        rounded corners, boxrule=0.5pt, left=2pt, right=2pt, top=2pt, bottom=2pt]
        \scriptsize
        \visualplaceholder[3cm]{Steam engine diagram}
        Entropy as ``energy unavailable for work.''
        Named from Greek \emph{entropia} (transformation).
      \end{tcolorbox}
    \end{column}

    \begin{column}{0.32\textwidth}
      % Column 2: Boltzmann
      \begin{tcolorbox}[colback=aiViolet!5, colframe=aiViolet!60,
        fonttitle=\bfseries\scriptsize, title={Boltzmann (1870s)},
        rounded corners, boxrule=0.5pt, left=2pt, right=2pt, top=2pt, bottom=2pt]
        \scriptsize
        \visualplaceholder[3cm]{Gas molecules in a box; Boltzmann's
          tombstone with $S = k \log W$}
        Entropy as ``number of arrangements.''
        \[S = k \log W\]
      \end{tcolorbox}
    \end{column}

    \begin{column}{0.32\textwidth}
      % Column 3: Shannon
      \begin{tcolorbox}[colback=aiViolet!5, colframe=aiViolet!60,
        fonttitle=\bfseries\scriptsize, title={Shannon (1948)},
        rounded corners, boxrule=0.5pt, left=2pt, right=2pt, top=2pt, bottom=2pt]
        \scriptsize
        \visualplaceholder[3cm]{Fair coin vs loaded coin; entropy
          formula diagram}
        Entropy as ``average surprise.''
        \[H = -\!\sum_i p_i \log p_i\]
        {\tiny $\sum$ = ``add up for every outcome $i$''}
      \end{tcolorbox}
    \end{column}
  \end{columns}

  \vspace{0.3em}
  {\small\centering\itshape The same formula. Three completely different domains.\par}

  \note{
    CROSS-REFERENCE: Shannon appears again! We saw him in Section~3
    (statistical language models, Slide~32). His 1948 paper is so
    foundational that it contributes to BOTH the probability and
    information theory sections. He modeled language statistically AND
    created the measure for how good that model is.

    Conceptual recipe for Shannon's formula: ``For each possible outcome,
    multiply its probability by the `surprise' (the logarithm of how
    unexpected it is). Add them all up. The result is the average
    surprise.''

    Anecdote: ``Shannon supposedly asked von Neumann what to call his new
    quantity. Von Neumann said: `Call it entropy. Nobody knows what
    entropy really is, so in a debate you will always have the
    advantage.'\,'' (May be apocryphal.)

    \verifytag{Clausius coined ``entropy'' in 1865 but the concept appears
    in his earlier 1850s work. Shannon--von Neumann anecdote is widely
    attributed but may be apocryphal.}
  }
\end{frame}

% ----------------------------------------------------------
% Slide 37: Cross-Entropy -- The Loss Function of LLMs
% ----------------------------------------------------------
\begin{frame}{Cross-Entropy --- THE Loss Function of LLMs}

  \begin{columns}[T]
    \begin{column}{0.46\textwidth}
      % Conceptual recipe
      \begin{tcolorbox}[colback=aiViolet!5, colframe=aiViolet!60,
        fonttitle=\bfseries\small, title={Recipe: Cross-Entropy Loss},
        rounded corners, boxrule=0.5pt, left=3pt, right=3pt, top=3pt, bottom=3pt]
        \footnotesize
        1. The true answer is one specific word.\\
        2. Look at the probability your model gave that word.\\
        3. Take the logarithm.\\
        4. Negate it.\\[3pt]
        High confidence in right answer \(\Rightarrow\) low loss.\\
        Low confidence in right answer \(\Rightarrow\) high loss.
      \end{tcolorbox}

      \vspace{0.3em}
      % Formal inset
      \begin{tcolorbox}[colback=gray!5, colframe=gray!40,
        fonttitle=\bfseries\small, title={Formal},
        rounded corners, boxrule=0.5pt, left=3pt, right=3pt, top=1pt, bottom=1pt]
        \footnotesize
        \[\text{Loss} = -\log\bigl(p_{\text{model}}(\text{correct token})\bigr)\]
        {\scriptsize Full form uses $\sum$ over all tokens, but since
        only one is correct, it simplifies to this.}
      \end{tcolorbox}
    \end{column}

    \begin{column}{0.51\textwidth}
      % Concrete example
      {\footnotesize ``The cat sat on the \rule{0.8cm}{0.4pt}''}

      \vspace{0.3em}
      \begin{tikzpicture}[scale=0.85]
        \begin{axis}[
          width=6.5cm, height=4cm,
          ybar,
          bar width=10pt,
          ylabel={\scriptsize Probability},
          ylabel style={at={(axis description cs:-0.05,0.5)}},
          symbolic x coords={mat, rug, floor, ground, ...},
          xtick=data,
          tick label style={font=\scriptsize},
          label style={font=\scriptsize},
          ymin=0, ymax=0.45,
          ytick={0,0.1,0.2,0.3,0.4},
          axis lines=left,
          nodes near coords={\scriptsize\pgfmathprintnumber\pgfplotspointmeta},
          every node near coord/.append style={font=\tiny},
        ]
          \addplot[fill=aiViolet!70] coordinates
            {(mat,0.30) (rug,0.25) (floor,0.20) (ground,0.15) (...,0.10)};
        \end{axis}
        \node[font=\scriptsize,aiViolet] at (1.2,-0.4)
          {Loss $= -\log(0.30) = 1.20$};
      \end{tikzpicture}

      \vspace{0.1em}
      {\scriptsize If model improves to 80\% for ``mat'':
       loss $= -\log(0.80) = 0.22$}

      \vspace{0.2em}
      \begin{insightbox}
        Training an LLM = minimizing cross-entropy = making the model
        more confident about the RIGHT next word.
      \end{insightbox}
    \end{column}
  \end{columns}

  \note{
    ``Every time an LLM gets a training example, it predicts the next
    token, checks how wrong it was using cross-entropy, then uses
    backpropagation (Section~2!) to improve. Cross-entropy is the bridge
    between information theory and gradient descent.''

    Bridge to earlier sections: Softmax (Section~3) produces the
    probabilities. Cross-entropy measures how good they are.
    Backpropagation (Section~2) improves them.
  }
\end{frame}

% ----------------------------------------------------------
% Slide 38: Surprisal and Perplexity
% ----------------------------------------------------------
\begin{frame}{Surprisal and Perplexity --- How We Evaluate LLMs}

  \begin{columns}[T]
    \begin{column}{0.47\textwidth}
      \begin{tcolorbox}[colback=aiViolet!5, colframe=aiViolet!60,
        fonttitle=\bfseries\small, title={Surprisal},
        rounded corners, boxrule=0.5pt, left=3pt, right=3pt, top=3pt, bottom=3pt]
        \footnotesize
        $\text{Surprisal} = -\log\bigl(p(\text{word})\bigr)$\\[3pt]
        High probability \(\Rightarrow\) low surprisal\\
        Low probability \(\Rightarrow\) high surprisal
      \end{tcolorbox}

      \vspace{0.3em}
      \begin{tcolorbox}[colback=aiViolet!5, colframe=aiViolet!60,
        fonttitle=\bfseries\small, title={Perplexity},
        rounded corners, boxrule=0.5pt, left=3pt, right=3pt, top=3pt, bottom=3pt]
        \footnotesize
        Average surprisal, exponentiated.\\[3pt]
        Lower perplexity = better model.\\[3pt]
        {\scriptsize ``GPT-4 has lower perplexity than GPT-3'' means
        it's less surprised by real text.}
      \end{tcolorbox}

      \vspace{0.3em}
      \begin{insightbox}
        A good language model is rarely surprised by real text.
        Shannon would have understood this immediately.
      \end{insightbox}
    \end{column}

    \begin{column}{0.5\textwidth}
      % Surprisal example
      {\footnotesize ``The cat sat on the \rule{0.8cm}{0.4pt}''}

      \vspace{0.4em}
      \begin{tikzpicture}[
        bar/.style={draw=none, rounded corners=1pt},
        lbl/.style={font=\scriptsize, anchor=west},
      ]
        % mat -- low surprisal
        \fill[aiViolet!30] (0, 2.4) rectangle (1.5, 2.8);
        \node[lbl] at (1.7, 2.6) {``mat'' --- low surprisal (expected)};

        % chandelier -- high surprisal
        \fill[aiViolet!60] (0, 1.4) rectangle (4.0, 1.8);
        \node[lbl] at (4.2, 1.6) {``chandelier''};
        \node[font=\tiny\itshape,gray,anchor=west] at (4.2,1.3) {high surprisal};

        % asdfghjkl -- very high
        \fill[aiViolet!90] (0, 0.4) rectangle (5.5, 0.8);
        \node[lbl] at (0.1, 0.1) {``asdfghjkl'' --- very high surprisal};

        % Arrow label
        \draw[->,thick,gray] (0,0.0) -- (5.8,0.0);
        \node[font=\tiny,gray,anchor=north] at (2.9,-0.1) {Surprisal \(\to\)};
      \end{tikzpicture}
    \end{column}
  \end{columns}

  \note{
    ``When researchers say GPT-4 has `lower perplexity' than GPT-3, they
    mean it's less surprised by real text --- it predicts better.
    Shannon would have understood this immediately.''

    Perplexity is the standard evaluation metric for language models. It
    connects directly to Shannon's information-theoretic framework.
  }
\end{frame}

% ----------------------------------------------------------
% Slide 39: KL Divergence [EXTENDED]
% ----------------------------------------------------------
\begin{frame}{KL Divergence --- Measuring the Gap Between Distributions}

  \begin{columns}[T]
    \begin{column}{0.47\textwidth}
      {\footnotesize
      \textbf{Solomon Kullback} (1907--1994) \&
      \textbf{Richard Leibler} (1914--2003), USA.\\
      Published KL divergence in 1951.}

      \vspace{0.3em}
      % Conceptual recipe
      \begin{tcolorbox}[colback=aiViolet!5, colframe=aiViolet!60,
        fonttitle=\bfseries\small, title={Recipe: KL Divergence},
        rounded corners, boxrule=0.5pt, left=3pt, right=3pt, top=2pt, bottom=2pt]
        \footnotesize
        For each possible outcome, ask:\\
        ``How much more surprising is this under model $Q$ compared to
        truth $P$?''\\[3pt]
        Average over all outcomes.\\[3pt]
        If $Q$ matches $P$ perfectly $\Rightarrow$ KL $= 0$.
      \end{tcolorbox}

      \vspace{0.3em}
      \begin{insightbox}
        In RLHF (the technique that makes ChatGPT helpful), KL divergence
        prevents the model from changing too much during fine-tuning.
      \end{insightbox}

      \vspace{0.2em}
      {\footnotesize\textcolor{gray}{[EXTENDED] --- can be cut without
      losing the narrative}}
    \end{column}

    \begin{column}{0.5\textwidth}
      % Two overlapping distributions
      \begin{tikzpicture}[scale=0.85]
        \begin{axis}[
          width=6.5cm, height=4.5cm,
          axis lines=left,
          xlabel={\scriptsize Outcome},
          ylabel={\scriptsize Probability},
          xtick=\empty, ytick=\empty,
          xmin=-3, xmax=4,
          ymin=0, ymax=0.5,
          tick label style={font=\tiny},
          label style={font=\tiny},
        ]
          % P (true)
          \addplot[aiViolet, very thick, smooth, domain=-3:4, samples=50]
            {0.4*exp(-x*x/2)};
          % Q (model)
          \addplot[accentOrange, very thick, smooth, dashed, domain=-3:4, samples=50]
            {0.35*exp(-(x-1)*(x-1)/2.5)};
          % Gap annotation
          \draw[<->, thick, red!60] (axis cs:0,0.35) -- (axis cs:1,0.28)
            node[midway, above, font=\tiny, text=red!70] {gap};
          \node[font=\scriptsize, aiViolet] at (axis cs:-1.5,0.35) {$P$ (true)};
          \node[font=\scriptsize, accentOrange] at (axis cs:2.5,0.3) {$Q$ (model)};
        \end{axis}
        \node[font=\scriptsize\itshape,gray] at (3.2,-0.3)
          {Bigger gap = worse model};
      \end{tikzpicture}

      \vspace{0.3em}
      % Formal
      \begin{tcolorbox}[colback=gray!5, colframe=gray!40,
        fonttitle=\bfseries\small, title={Formal},
        rounded corners, boxrule=0.5pt, left=3pt, right=3pt, top=1pt, bottom=1pt]
        \footnotesize
        $\displaystyle D_{\text{KL}}(P \| Q) = \sum_i P(i)\,\log\frac{P(i)}{Q(i)}$
      \end{tcolorbox}
    \end{column}
  \end{columns}

  \note{
    This slide is marked [EXTENDED] and can be cut if time is short.

    ``In RLHF (Reinforcement Learning from Human Feedback), KL divergence
    is used to prevent the model from diverging too far from its base
    behavior during fine-tuning.''

    \verifytag{Kullback--Leibler 1951 paper ``On Information and
    Sufficiency'' --- confirm title and year.}
  }
\end{frame}

% ----------------------------------------------------------
% Slide 40: Tokenization -- Carving Language into Pieces
% ----------------------------------------------------------
\begin{frame}{Tokenization --- Carving Language into Pieces}

  \begin{columns}[T]
    \begin{column}{0.46\textwidth}
      {\footnotesize
      \textbf{George Kingsley Zipf} (1902--1950, USA).\\
      \alert{Zipf's law} (1935/1949): word frequency is inversely
      proportional to its rank. ``The'' is \#1, ``of'' is \#2, \ldots}

      \vspace{0.3em}
      {\footnotesize
      \textbf{Rico Sennrich}, \textbf{Barry Haddow},
      \textbf{Alexandra Birch} (Edinburgh, 2016).\\
      Adapted \alert{Byte Pair Encoding} (BPE) for neural machine
      translation. BPE is the tokenization method used by GPT and most
      modern LLMs.}

      \vspace{0.3em}
      \begin{insightbox}
        The AI doesn't see words --- it sees tokens. BPE gives short
        tokens to common patterns (Zipf's law!) and splits rare words
        into pieces.
      \end{insightbox}
    \end{column}

    \begin{column}{0.51\textwidth}
      % Tokenization example
      \begin{tcolorbox}[colback=aiViolet!5, colframe=aiViolet!60,
        fonttitle=\bfseries\small, title={BPE Tokenization Example},
        rounded corners, boxrule=0.5pt, left=3pt, right=3pt, top=3pt, bottom=3pt]
        \footnotesize
        \textbf{Input:}\\
        {\ttfamily The mathematician calculated the eigenvalue}\\[4pt]
        \textbf{Tokens:}\\
        \colorbox{aiViolet!15}{\ttfamily The}\;
        \colorbox{aiViolet!25}{\ttfamily\, mathematic}\,%
        \colorbox{aiViolet!35}{\ttfamily ian}\;
        \colorbox{aiViolet!15}{\ttfamily\, calculated}\;
        \colorbox{aiViolet!15}{\ttfamily\, the}\;
        \colorbox{aiViolet!25}{\ttfamily\, eigen}\,%
        \colorbox{aiViolet!35}{\ttfamily value}\\[4pt]
        {\scriptsize Common words \(\to\) single token\\
        Rare words \(\to\) split into pieces}
      \end{tcolorbox}

      \vspace{0.3em}
      % Zipf's law sketch
      \begin{tikzpicture}[scale=0.8]
        \begin{axis}[
          width=5.5cm, height=3cm,
          axis lines=left,
          xlabel={\tiny Word rank},
          ylabel={\tiny Frequency},
          xtick={1,10,100,1000},
          xmode=log,
          ymode=log,
          ymin=1, ymax=100000,
          tick label style={font=\tiny},
          label style={font=\tiny},
        ]
          \addplot[aiViolet, thick, domain=1:1000, samples=30] {80000/x};
          \node[font=\tiny] at (axis cs:3,40000) {``the''};
          \node[font=\tiny] at (axis cs:50,800) {``eigenvalue''};
        \end{axis}
      \end{tikzpicture}
    \end{column}
  \end{columns}

  \note{
    ``Common words like `the' get one token. Uncommon words like
    `eigenvalue' get split into `eigen' + `value.' This encoding scheme
    is deeply connected to information theory: give short codes to
    frequent patterns.''

    Historical lineage: Philip Gage published BPE as a compression
    algorithm in 1994. Sennrich et al.\ adapted it for NLP in 2016.

    \verifytag{Zipf's law appears in his 1935 book and more fully in his
    1949 book. Sennrich et al.\ 2016 ACL paper --- confirm venue.}
  }
\end{frame}

% ----------------------------------------------------------
% Slide 41: Convergence Diagram Update -- Information Theory
% ----------------------------------------------------------
\begin{frame}{Building the Transformer --- Information Theory}

  \begin{columns}[T]
    \begin{column}{0.45\textwidth}
      % Convergence diagram placeholder
      \visualplaceholder[5cm]{Convergence diagram: Sections 1--3 lit
        (Linear Algebra, Calculus, Probability). Section~4
        ``Information Theory'' now lights up in red.
        Label: ``Cross-entropy = the loss function.''}
    \end{column}

    \begin{column}{0.52\textwidth}
      \begin{tcolorbox}[colback=aiViolet!5, colframe=aiViolet,
        fonttitle=\bfseries, title={Section 4 Summary},
        rounded corners, boxrule=1pt, left=4pt, right=4pt, top=3pt, bottom=3pt]
        \footnotesize
        \begin{tabular}{@{}ll@{}}
          \toprule
          \textbf{Who / When} & \textbf{Contribution}\\
          \midrule
          Clausius (1865) & Thermodynamic entropy\\
          Boltzmann (1870s) & Statistical entropy, $S = k\log W$\\
          Shannon (1948) & Information entropy, the ``bit''\\
          Kullback--Leibler (1951) & Divergence between distributions\\
          Zipf (1935/1949) & Word frequency laws\\
          Sennrich et al.\ (2016) & BPE tokenization\\
          \bottomrule
        \end{tabular}
      \end{tcolorbox}

      \vspace{0.3em}
      \begin{insightbox}
        LLMs are trained by minimizing cross-entropy loss.
        Shannon's 1948 formula is literally the objective function.
      \end{insightbox}
    \end{column}
  \end{columns}

  \note{
    Connection arrow: ``LLMs are trained by minimizing cross-entropy
    loss. Shannon's 1948 formula is literally the objective function.''

    Transition: ``We have the math for learning and the math for
    measuring. But we need something to run it on. Enter: logic and
    computation.''
  }
\end{frame}


% ************************************************************
% SECTION 5: LOGIC & COMPUTATION (Slides 42-48)
% "The Substrate"
% Section color: aiViolet
%
% NARRATIVE TENSION: Logicians dreamed of mechanical reasoning.
% They proved it was impossible. Then neural networks did it
% anyway.
% ************************************************************
\section{Logic \& Computation}

% ----------------------------------------------------------
% Slide 42: Section Divider
% ----------------------------------------------------------
{
\setbeamercolor{background canvas}{bg=aiViolet}
\begin{frame}[plain]
  \centering
  \vspace{2cm}
  \textcolor{white}{\Huge\bfseries Logic \& Computation}\\[0.5em]
  \textcolor{white}{\Huge\bfseries The Substrate}\\[1em]
  \textcolor{white!80}{\large 1854--1970s}\\[1.5em]
  \textcolor{white!60}{\normalsize Boole \(\cdot\) Frege \(\cdot\) G\"odel \(\cdot\) Turing \(\cdot\) von Neumann}\\[1em]
  \textcolor{white!70}{\itshape\normalsize ``Logicians dreamed of mechanical thought.\\
    They proved it was impossible.\\
    Then neural networks did it anyway.''}

  \note{
    Transition: ``We have the math for learning and the math for
    measuring. But we need something to run it on.''

    This section is structured around a paradox: logicians tried to
    reduce all reasoning to mechanical rules, Turing and G\"odel proved
    that's impossible, and yet neural networks learned to do things
    formal logic cannot: translate, write poetry, reason by analogy.

    The machines that logic built transcended the limitations that logic
    discovered.

    This section covers \(\sim\)8 minutes.
  }
\end{frame}
}

% ----------------------------------------------------------
% Slide 43: Boolean Algebra -- The Language of Computers
% ----------------------------------------------------------
\begin{frame}{Boolean Algebra --- The Language of Computers (1854)}

  \begin{columns}[T]
    \begin{column}{0.45\textwidth}
      \begin{personbox}[George Boole (1815--1864)]{aiViolet}
        \textbf{Origin:} Lincoln, England; worked at Queen's College
        Cork, Ireland\\[3pt]
        Published \emph{An Investigation of the Laws of Thought} (1854).
        Self-taught mathematician who showed that \textbf{logical
        reasoning} could be expressed as \textbf{algebra}: variables
        take values 0 (false) or 1 (true).
      \end{personbox}

      \vspace{0.3em}
      \begin{insightbox}
        Boole was a self-taught son of a shoemaker who became a
        university professor. His algebra is now the foundation of every
        digital device on Earth.
      \end{insightbox}
    \end{column}

    \begin{column}{0.52\textwidth}
      % Boolean logic gates
      \begin{tikzpicture}[
        gate/.style={draw=aiViolet, thick, rounded corners, fill=aiViolet!8,
          minimum width=1.6cm, minimum height=0.7cm, font=\small\bfseries,
          align=center},
        tbl/.style={font=\tiny, inner sep=2pt},
        >=Stealth
      ]
        % AND gate
        \node[gate] (and) at (0, 2.5) {AND};
        \node[tbl,right=0.3cm of and] {%
          \begin{tabular}{cc|c}
            A & B & Out\\ \hline
            0 & 0 & 0\\
            0 & 1 & 0\\
            1 & 0 & 0\\
            1 & 1 & 1
          \end{tabular}
        };

        % OR gate
        \node[gate] (or) at (0, 0.8) {OR};
        \node[tbl,right=0.3cm of or] {%
          \begin{tabular}{cc|c}
            A & B & Out\\ \hline
            0 & 0 & 0\\
            0 & 1 & 1\\
            1 & 0 & 1\\
            1 & 1 & 1
          \end{tabular}
        };

        % NOT gate
        \node[gate] (not) at (0, -0.5) {NOT};
        \node[tbl,right=0.3cm of not] {%
          \begin{tabular}{c|c}
            A & Out\\ \hline
            0 & 1\\
            1 & 0
          \end{tabular}
        };
      \end{tikzpicture}

      \vspace{0.3em}
      {\small\centering Every GPU in every AI data center is made of
      \emph{billions} of these gates.\par}
    \end{column}
  \end{columns}

  \note{
    ``Boolean algebra reduced logical reasoning to symbols and rules.
    Every digital computer --- from your phone to a data center full of
    GPUs --- is built from AND, OR, and NOT gates. Billions of them.''

    Boole's work was abstract and philosophical. He had no idea it would
    become the foundation of hardware.
  }
\end{frame}

% ----------------------------------------------------------
% Slide 44: The Dream and Its Limits -- Frege to Godel to Turing
% ----------------------------------------------------------
\begin{frame}{The Dream and Its Limits --- From Frege to Turing}

  {\small After Boole, logicians dreamed of reducing ALL reasoning to
  mechanical symbol manipulation. They almost succeeded --- and then
  proved it was \textbf{impossible}.}

  \vspace{0.3em}
  % Timeline
  \begin{tikzpicture}[
    person/.style={draw=aiViolet!60, rounded corners, fill=aiViolet!5,
      font=\scriptsize, inner sep=3pt, text width=3.1cm, align=left,
      minimum height=1.6cm},
    arr/.style={-Stealth, thick, aiViolet!50},
  ]
    \node[person] (frege) at (0, 0) {%
      \textbf{Frege} (1879)\\
      First formal logical system.\\
      A precursor to programming languages.};
    \node[person] (rw) at (3.8, 0) {%
      \textbf{Russell \& Whitehead}\\
      \emph{Principia} (1910--13)\\
      362 pages to prove $1{+}1{=}2$.};
    \node[person] (godel) at (7.6, 0) {%
      \textbf{G\"odel} (1931)\\
      \alert{Incompleteness:} true statements exist that \emph{cannot}
      be proved.};
    \node[person] (turing) at (11.4, 0) {%
      \textbf{Turing} (1936)\\
      \alert{Halting problem:} no algorithm can decide if any program
      will stop.};

    \draw[arr] (frege.east) -- (rw.west);
    \draw[arr] (rw.east) -- (godel.west);
    \draw[arr] (godel.east) -- (turing.west);
  \end{tikzpicture}

  \vspace{0.3em}
  \begin{insightbox}
    The irony: neural networks, built from Boole's logic gates, learned
    to translate, reason, and create --- not by following rules, but by
    learning patterns from data. The machines that logic built went
    beyond what logic can do.
  \end{insightbox}

  \note{
    Key persons (brief):
    \begin{itemize}\setlength\itemsep{1pt}
      \item \textbf{Gottlob Frege} (1848--1925, Wismar, Germany).
        Created the first formal logic powerful enough to express math
        (\emph{Begriffsschrift}, 1879).
      \item \textbf{Bertrand Russell} (1872--1970) \&
        \textbf{Alfred North Whitehead} (1861--1947).
        \emph{Principia Mathematica} (1910--1913): derive ALL math from
        logical axioms.
      \item \textbf{Kurt G\"odel} (1906--1978, Brno / Princeton).
        Incompleteness theorems (1931).
      \item \textbf{Alan Turing} (1912--1954, London). ``On Computable
        Numbers'' (1936): the Turing machine and the halting problem.
    \end{itemize}

    ``They tried to reduce all reasoning to logic. G\"odel proved some
    truths escape any formal system. Turing proved some questions can't
    be answered by any algorithm. And yet --- neural networks, built from
    Boole's logic gates, learned to translate, reason, and create.''

    \verifytag{Page number 362 for 1+1=2 in Principia is commonly cited
    but may be approximate.}
  }
\end{frame}

% ----------------------------------------------------------
% Slide 45: Von Neumann Architecture [EXTENDED]
% ----------------------------------------------------------
\begin{frame}{Von Neumann Architecture --- The Hardware Blueprint (1945)}

  \begin{columns}[T]
    \begin{column}{0.45\textwidth}
      \begin{personbox}[John von Neumann (1903--1957)]{aiViolet}
        \textbf{Born:} Neumann J\'anos Lajos\\
        \textbf{Origin:} Budapest, Hungary; worked at Princeton, USA\\[3pt]
        Published \emph{First Draft of a Report on the EDVAC} (1945),
        describing the \alert{stored-program architecture}: instructions
        and data share the same memory.
      \end{personbox}

      \vspace{0.3em}
      {\footnotesize\textcolor{gray}{[EXTENDED] --- can be cut;
      summarize verbally}}

      \vspace{0.2em}
      {\small\itshape LLM inference runs on this architecture. But
      training requires massively parallel hardware. Enter the GPU.}
    \end{column}

    \begin{column}{0.52\textwidth}
      % Von Neumann architecture diagram
      \begin{tikzpicture}[
        comp/.style={draw=aiViolet, thick, rounded corners=4pt,
          fill=aiViolet!10, minimum width=2.2cm, minimum height=0.9cm,
          font=\small\bfseries, align=center},
        bus/.style={draw=aiViolet, thick, -Stealth, line width=1.5pt},
        dbus/.style={draw=aiViolet, thick, Stealth-Stealth, line width=1.5pt},
      ]
        \node[comp] (cpu) at (0, 2) {CPU\\{\scriptsize Control + ALU}};
        \node[comp] (mem) at (0, 0) {Memory\\{\scriptsize Programs + Data}};
        \node[comp] (inp) at (-2.5, 0) {Input};
        \node[comp] (out) at (2.5, 0) {Output};

        \draw[dbus] (cpu) -- (mem);
        \draw[bus] (inp) -- (mem);
        \draw[bus] (mem) -- (out);
      \end{tikzpicture}

      \vspace{0.4em}
      {\small\centering This basic design has powered nearly every
      computer for 80~years.\par}
    \end{column}
  \end{columns}

  \note{
    ``The stored-program concept: keep both the program and its data in
    the same memory. Nearly every computer follows this design.''

    Note: the 1945 draft was controversial --- it was circulated under
    von Neumann's name alone, but was based on collaborative work by
    Eckert and Mauchly. We present the architecture concept, not the
    priority dispute.

    \verifytag{First Draft of a Report on the EDVAC, 1945 --- confirm
    attribution context.}
  }
\end{frame}

% ----------------------------------------------------------
% Slide 46: Floating Point Arithmetic
% ----------------------------------------------------------
\begin{frame}{Floating Point --- How Computers Handle Real Numbers}

  \begin{columns}[T]
    \begin{column}{0.48\textwidth}
      {\footnotesize
      \textbf{Konrad Zuse} (1910--1995, Berlin, Germany).\\
      Built the \alert{Z3} (1941), the first programmable computer using
      floating-point arithmetic.}

      \vspace{0.3em}
      % Floating-point layout
      \begin{tcolorbox}[colback=aiViolet!5, colframe=aiViolet!60,
        fonttitle=\bfseries\small, title={How 3.14159 is stored},
        rounded corners, boxrule=0.5pt, left=3pt, right=3pt, top=3pt, bottom=3pt]
        \footnotesize
        \begin{tikzpicture}
          \fill[aiViolet!20] (0,0) rectangle (0.5,0.5);
          \node[font=\tiny] at (0.25,0.25) {sign};
          \fill[aiViolet!35] (0.6,0) rectangle (2.6,0.5);
          \node[font=\tiny] at (1.6,0.25) {exponent (8 bits)};
          \fill[aiViolet!50] (2.7,0) rectangle (5.8,0.5);
          \node[font=\tiny,white] at (4.25,0.25) {mantissa (23 bits)};
        \end{tikzpicture}

        \vspace{2pt}
        {\scriptsize IEEE 754 standard (1985)}
      \end{tcolorbox}

      \vspace{0.3em}
      \begin{insightbox}
        Every weight in an LLM is a floating-point number. Large models
        have hundreds of billions of them.
      \end{insightbox}
    \end{column}

    \begin{column}{0.49\textwidth}
      % Precision comparison
      \begin{tcolorbox}[colback=gray!5, colframe=gray!50,
        fonttitle=\bfseries\small,
        title={Precision vs.\ Speed Trade-off},
        rounded corners, boxrule=0.5pt, left=3pt, right=3pt, top=3pt, bottom=3pt]
        \footnotesize
        \begin{tabular}{@{}lcc@{}}
          \toprule
          \textbf{Format} & \textbf{Bits} & \textbf{Use}\\
          \midrule
          FP32 & 32 & Traditional training\\
          FP16 / BF16 & 16 & Modern LLM training\\
          INT8 & 8 & Fast inference\\
          INT4 & 4 & Aggressive compression\\
          \bottomrule
        \end{tabular}

        \vspace{4pt}
        {\scriptsize Fewer bits = less precise but faster and smaller.\\
        Modern LLMs accept small rounding errors in exchange for fitting
        bigger models.}
      \end{tcolorbox}

      \vspace{0.3em}
      {\small\centering\itshape From Zuse's Z3 (1941) to today's
      GPU clusters: the same fundamental idea, vastly more scale.\par}
    \end{column}
  \end{columns}

  \note{
    ``Modern LLMs use reduced-precision floating point (16-bit or even
    8-bit instead of 32-bit) to save memory and speed up computation.
    This means accepting small rounding errors in exchange for being able
    to fit a bigger model.''

    \verifytag{Z3 as first floating-point computer is standard but note
    it was electromechanical, not electronic.}
  }
\end{frame}

% ----------------------------------------------------------
% Slide 47: GPU Computing -- Why Matrix Multiplication Won
% ----------------------------------------------------------
\begin{frame}{GPU Computing --- Why Matrix Multiplication Won}

  \begin{columns}[T]
    \begin{column}{0.47\textwidth}
      {\footnotesize
      \textbf{NVIDIA} released \alert{CUDA} (2007): general-purpose
      computing on GPUs.}

      \vspace{0.2em}
      {\footnotesize
      \textbf{Alex Krizhevsky}, \textbf{Ilya Sutskever},
      \textbf{Geoffrey Hinton} --- \alert{AlexNet} (2012): GPUs could
      train deep neural networks dramatically faster, launching the deep
      learning era.}

      \vspace{0.3em}
      \begin{insightbox}
        Without GPUs, deep learning would still be a curiosity.
        Video game hardware accidentally enabled AI.
      \end{insightbox}

      \vspace{0.3em}
      {\small\itshape ``AlexNet showed GPUs could train networks
      10--100$\times$ faster than CPUs. That single result changed the
      trajectory of AI.''}
    \end{column}

    \begin{column}{0.5\textwidth}
      \visualplaceholder[5cm]{CPU vs GPU comparison:\\CPU: 4--16 powerful cores\\GPU: 5,000--16,000 simple cores\\Neural networks = parallel matrix math = perfect for GPU}

      \vspace{0.3em}
      {\small\centering Neural network training: millions of
      \emph{independent} matrix multiplications\\
      \(\Rightarrow\) perfect for GPU.\par}
    \end{column}
  \end{columns}

  \note{
    ``Without GPUs, deep learning would still be a curiosity. The
    AlexNet result in 2012 showed that GPUs could train networks
    10--100$\times$ faster than CPUs. That single result changed the
    trajectory of AI.''

    Key insight: GPUs were designed for rendering video game graphics ---
    which is also parallel matrix math. The fit with neural networks was
    accidental and transformative.

    \verifytag{AlexNet: Krizhevsky et al., NeurIPS (then NIPS) 2012 ---
    confirm venue.}
  }
\end{frame}

% ----------------------------------------------------------
% Slide 48: Convergence Diagram Update -- Logic & Computation
% ----------------------------------------------------------
\begin{frame}{Building the Transformer --- Logic \& Computation}

  \begin{columns}[T]
    \begin{column}{0.45\textwidth}
      % Convergence diagram placeholder
      \visualplaceholder[5cm]{Convergence diagram: Sections 1--4 lit
        (Linear Algebra, Calculus, Probability, Information Theory).
        Section~5 ``Logic \& Computation'' now lights up in gray/steel.
        Label: ``GPUs = the hardware.''}

      \vspace{0.3em}
      \begin{insightbox}
        Boolean logic built the hardware. Turing defined computation.
        GPUs made deep learning possible. And neural networks surpassed
        what formal logic can do.
      \end{insightbox}
    \end{column}

    \begin{column}{0.52\textwidth}
      \begin{tcolorbox}[colback=aiViolet!5, colframe=aiViolet,
        fonttitle=\bfseries, title={Section 5 Summary},
        rounded corners, boxrule=1pt, left=4pt, right=4pt, top=3pt, bottom=3pt]
        \footnotesize
        \begin{tabular}{@{}ll@{}}
          \toprule
          \textbf{Who / When} & \textbf{Contribution}\\
          \midrule
          Boole (1854) & Boolean algebra\\
          Frege (1879) & Formal logic\\
          Russell/Whitehead (1910--13) & \emph{Principia Mathematica}\\
          G\"odel (1931) & Incompleteness theorems\\
          Turing (1936) & Computability, Turing machine\\
          Von Neumann (1945) & Stored-program architecture\\
          Zuse (1941) & Floating-point computation\\
          CUDA (2007) + AlexNet (2012) & GPU revolution\\
          \bottomrule
        \end{tabular}
      \end{tcolorbox}
    \end{column}
  \end{columns}

  \note{
    Connection arrow: ``Boolean logic built the hardware. Turing defined
    computation. GPUs made deep learning possible. And neural networks
    surpassed what formal logic can do.''

    Transition: ``We have math AND hardware. But here's a deep question:
    WHY should a pile of matrix multiplications be able to learn
    ANYTHING? For the answer, we need the theory of function
    approximation.''

    This marks the natural end of Session~2 (Slides 26--48,
    \(\sim\)26 minutes). Natural break line: ``We have the math and the
    hardware. But can we actually build a brain?''
  }
\end{frame}
